\subsection{Negligible positive functions}

DEFINITION 1. --- Given a measure \(\mu\) on a locally compact space \(X\) , a numerical function \(f \geqslant 0\) (finite or not) defined on \(X\) is said to be negligible for the measure \(\mu\) if \(|\mu|^*(f) = 0\) .

We then also say that f is \(\mu\) -negligible, or simply negligible if no confusion can result.

PROPOSITION 1. --- If \(f\) is a negligible function \(\geqslant 0\) , then every numerical function \(g\) such that \(0 \leqslant g \leqslant \alpha f\) ( \(\alpha\) a scalar \(> 0\) ) is negligible.

For, \(0 \leqslant |\mu|^*(g) \leqslant \alpha |\mu|^*(f) = 0\)

PROPOSITION 2. --- The sum and upper envelope of a sequence \((f_{n})\) of negligible functions \(\geqslant0\) are negligible.

For, \(|\mu|^*(\sum_n f_n) \leqslant \sum_n |\mu|^*(f_n) = 0\) (§1, No.~3, Prop. 13) and \(\sup_n f_n \leqslant \sum_n f_n\) .

PROPOSITION 3. --- For a lower semi-continuous function \(f \geqslant 0\) on \(X\) to be negligible, it is necessary and sufficient that \(f\) be zero on the support of \(\mu\) .

If \(|\mu|^*(f) = 0\) then \(|\mu|(g) = 0\) for every function \(g \in \mathcal{H}_+\) such that \(g \leqslant f\) ; it follows (Ch. III, §2, No.~3, Prop. 9) that \(g\) is zero on the support S of \(\mu\) ; since \(f\) is the upper envelope of the functions \(g \in \mathcal{H}_+\) such that \(g \leqslant f\) (§1, No.~1, Lemma), \(f(x) = 0\) on S. Conversely, if \(f(x) = 0\) on S then \(g(x) = 0\) on S for every function \(g \in \mathcal{H}_+\) such that \(g \leqslant f\) , therefore (Ch. III, §2, No.~3, Prop. 8) \(|\mu|(g) = 0\) , which, by definition, implies that \(|\mu|^*(f) = 0\) .

